\documentclass[article]{jss}

%% -- LaTeX packages and custom commands ---------------------------------------

\usepackage{lmodern}
\usepackage{amsmath}

\newcommand{\class}[1]{`\code{#1}'}
\newcommand{\fct}[1]{\code{#1()}}

%% -- Article metainformation --------------------------------------------------

\author{
  Phan Nguyen Huong Le\\ Rutgers University \And MengYang Yi\\ Johns Hopkins University \AND Philip He\\ Rutgers University, Daiichi Sankyo Inc.
}
\Plainauthor{Phan Nguyen Huong Le, MengYang Yi, Philip He}

\title{Interactive Visualization of Graphical Multiple Testing Procedures: An \proglang{R} \pkg{shiny} Interface for Clinical Trial Design}
\Plaintitle{Interactive Visualization of Graphical Multiple Testing Procedures: An R Shiny Interface for Clinical Trial Design}
\Shorttitle{\pkg{graphMTP}: Interactive Visualization of Graphical Multiple Testing Procedures}

\Abstract{
Graphical multiple testing procedures are widely used in confirmatory clinical trials to control the family-wise error rate across multiple endpoints, treatment comparisons, and patient populations. Although the graphical framework unifies classical procedures such as Bonferroni, Holm, fixed-sequence, fallback, and gatekeeping strategies, specifying valid graphs and tracking significance-level reallocation after each rejection step can be computationally and conceptually
difficult. These challenges become more pronounced in group sequential settings, where hypotheses are tested repeatedly across interim analyses. We present \pkg{graphMTP}, an interactive web-based interface implemented in \proglang{R} with \pkg{shiny}, \pkg{TrialSimulator}, and \pkg{visNetwork},
for the design, visualization, and exploration of graphical multiple testing procedures. The software allows users to construct hypothesis graphs interactively, assign initial weights and transition weights with real-time constraint checking, and simulate rejection paths to visualize $\alpha$-propagation dynamically. By combining a visual interface with a statistically validated backend, \pkg{graphMTP} supports both the specification of complex testing strategies and their communication to interdisciplinary clinical trial teams. Fixed-design and group sequential examples illustrate how interactive visualization can improve workflow, interpretation, and reproducibility in multiplicity planning.}

\Keywords{family-wise error rate, graphical multiple testing, group sequential design, \proglang{R}, \pkg{shiny}, \pkg{TrialSimulator}, \pkg{visNetwork}}
\Plainkeywords{family-wise error rate, graphical multiple testing, group sequential design, R, shiny, TrialSimulator, visNetwork}

\Address{
  Phan Nguyen Huong Le\\
  Department of Statistics\\
  Rutgers University\\
  Piscataway, New Jersey, USA\\
  E-mail: \email{hpl14@scarletmail.rutgers.edu}
}

\Address{
  MengYang Yi\\
  Department of Biostatistics\\
  Johns Hopkins University\\
  Baltimore, Maryland, USA\\
  E-mail: \email{your-jhu-email@jhu.edu}
}

\Address{
  Philip He\\
  Department of Statistics\\
  Rutgers University\\
  Piscataway, New Jersey, USA\\
  \emph{and}\\
  Daiichi Sankyo Inc.\\
  Basking Ridge, New Jersey, USA\\
  E-mail: \email{your-email-here}
}

\begin{document}

\section{Introduction}

Confirmatory clinical trials frequently involve the simultaneous evaluation of
multiple hypotheses, including multiple endpoints, treatment comparisons, dose
groups, and patient subpopulations. When several hypotheses are tested within a
single study, the probability of making at least one false positive conclusion
can exceed the nominal significance level unless multiplicity is handled
appropriately \citep{FDA2017,EMA2017}. Strong control of the family-wise error
rate (FWER) is therefore a central requirement in the design and analysis of
confirmatory trials \citep{FDA2017,EMA2017}.

Classical multiplicity-adjustment methods include procedures such as the
Bonferroni correction, Holm's step-down procedure \citep{Holm1979}, and
fixed-sequence testing \citep{Wiens2003}. While these approaches provide strong
control of the FWER, they can be restrictive when complex hierarchical
relationships exist among hypotheses. Modern clinical trials often involve
structured testing strategies, such as gatekeeping procedures
\citep{Dmitrienko2003}, fallback strategies \citep{Wiens2003}, or sequential
testing across endpoints and patient populations. These settings require more
flexible methods that can represent and manage dependencies among hypotheses
while maintaining statistical validity.

Graphical multiple testing procedures provide such a flexible and unified
framework \citep{Bretz2009,Bretz2011}. In this framework, hypotheses are
represented as nodes in a directed graph, and weighted edges determine how the
significance level $\alpha$ is redistributed when a hypothesis is rejected.
Each node is initially assigned a portion of the global significance level, and
the resulting procedure can be interpreted through the closed testing principle
of \citet{Marcus1976}. This representation unifies a broad class of multiple
testing procedures while preserving strong FWER control
\citep{Bretz2009,Bretz2011}. The framework has also been extended to group
sequential designs with interim analyses, where alpha-spending functions
govern the allocation of significance level across both hypotheses and time
points \citep{LanDeMets1983,ProschanLanWittes2006,MaurerBretz2013}.

Despite their methodological appeal, graphical procedures can be difficult to
specify, validate, and communicate in practice. Users must ensure that initial
alpha allocations and outgoing edge weights satisfy the required constraints,
and they must understand how significance level propagates across the graph
after each rejection step. These tasks are manageable for small examples, but
they become substantially more difficult as the testing strategy grows more
complex, especially when interim analyses and group sequential boundaries are
incorporated.

Existing software supports important parts of this workflow, but not the full
interactive workflow considered here. The \pkg{gMCP} package implements
graphical multiple testing procedures in \proglang{R} and has provided an
earlier graphical user interface for Bonferroni-based graphical procedures
\citep{Bretz2011,gMCP2024}. The more recent \pkg{graphicalMCP} package offers a
code-based implementation of graph-based multiple comparison procedures
\citep{graphicalMCP2024}. However, neither package natively supports the group
sequential extension of \citet{MaurerBretz2013}, which requires additional
integration with group sequential design tools such as \pkg{gsDesign}
\citep{gsDesign2024}. As a result, there remains a need for a lightweight,
browser-based interface that combines interactive graph construction, real-time
validation, and dynamic visualization of rejection paths within a unified
workflow.

We address this need with \pkg{graphMTP}, an interactive \proglang{R}
application built with \pkg{shiny} \citep{shiny2024} for constructing and
exploring graphical multiple testing procedures in clinical trials. The
application allows users to create hypothesis graphs visually, assign initial
alpha weights, define transition rules, and simulate rejection sequences with
immediate visual feedback. In addition, \pkg{graphMTP} integrates fixed-design
graphical procedures with group sequential functionality through
\pkg{TrialSimulator} \citep{TrialSimulator2024} and uses
\pkg{visNetwork} \citep{visNetwork2023} for interactive graph visualization.
The software is intended both as a practical design aid for statisticians and
as a communication tool for interdisciplinary clinical trial teams.

The contribution of this paper is twofold. First, we describe the design and
implementation of \pkg{graphMTP} as open-source statistical software for the
interactive specification and exploration of graphical multiple testing
procedures. Second, we demonstrate through reproducible examples how the
software can improve workflow, interpretation, and communication in multiplicity
planning for clinical trials.

The remainder of the paper is organized as follows.
Section~\ref{sec:framework} reviews the statistical framework underlying
graphical multiple testing procedures. Section~\ref{sec:software} describes the
design and scope of \pkg{graphMTP}. Section~\ref{sec:implementation} presents
the application architecture and implementation details.
Section~\ref{sec:workflow} describes the user workflow and interface.
Section~\ref{sec:examples} provides reproducible examples in fixed-design and
group sequential settings. Section~\ref{sec:comparison} compares
\pkg{graphMTP} with related software and discusses current limitations and
future extensions. Section~\ref{sec:availability} concludes with software
availability and reproducibility details.

%--------------- STATS FRAMEWORK -----------------

\section{Statistical framework}
\label{sec:framework}

This section briefly summarizes the statistical framework underlying
\pkg{graphMTP}. We follow the notation and formulation of
\citet{Marcus1976}, \citet{Bretz2009}, \citet{Bretz2011}, and
\citet{MaurerBretz2013}.

\subsection{Family-wise error rate and closed testing}

Consider the problem of testing $m$ elementary null hypotheses
$H_1, \ldots, H_m$, and let
\[
I = \{1, \ldots, m\}
\]
denote the associated index set. In confirmatory clinical trials, the
family-wise error rate (FWER), defined as the probability of rejecting at least
one true null hypothesis, is typically required to be controlled in the strong
sense at a pre-specified level $\alpha$.

A general principle for constructing such procedures is the closed testing
principle of \citet{Marcus1976}. For every nonempty subset $J \subseteq I$,
define the intersection hypothesis
\[
H_J = \bigcap_{j \in J} H_j.
\]
The corresponding closed testing procedure rejects an elementary hypothesis
$H_i$, $i \in I$, if and only if all intersection hypotheses $H_J$ with
$i \in J \subseteq I$ are rejected by their corresponding level-$\alpha$ tests.
By construction, the resulting procedure controls the FWER in the strong sense
at level $\alpha$. However, direct implementation becomes cumbersome as the
number of hypotheses increases, because the closure contains $2^m - 1$
nonempty intersection hypotheses.

\subsection{Graphical multiple testing procedures}

An important class of closed testing procedures is obtained by applying
weighted Bonferroni tests to the intersection hypotheses
\citep{Bretz2009,Bretz2011}. For each nonempty subset $J \subseteq I$, assume
a collection of weights $w_j(J)$, $j \in J$, satisfying
\[
0 \leq w_j(J) \leq 1
\qquad \text{and} \qquad
\sum_{j \in J} w_j(J) \leq 1.
\]
The weighted Bonferroni test rejects $H_J$ if
\[
p_j \leq w_j(J)\alpha
\]
for at least one $j \in J$, where $p_j$ denotes the unadjusted $p$-value for
$H_j$.

Graphical multiple testing procedures provide a convenient representation of
such weighting strategies. Following \citet{Bretz2011}, a graphical weighting
strategy is specified through the weights $w_i(I)$, $i \in I$, for the global
intersection hypothesis $H_I$, together with a transition matrix
\[
G = (g_{ij}),
\]
where
\[
0 \leq g_{ij} \leq 1, \qquad g_{ii} = 0,
\qquad \text{and} \qquad \sum_{j=1}^m g_{ij} \leq 1
\]
for all $i,j \in I$. Here, $w_i(I)$ denotes the initial weight assigned to
$H_i$, and $g_{ij}$ indicates the fraction of the local significance level at
$H_i$ that is propagated to $H_j$ if $H_i$ is rejected.

The full collection of weights $w_j(J)$ for the closure is induced from
$w_i(I)$ and $G$ through the graph-updating algorithm described by
\citet{Bretz2011}. Many standard multiple testing procedures, including
weighted Bonferroni, Holm, fixed-sequence, fallback, and gatekeeping
procedures, can be expressed within this framework
\citep{Bretz2009,Bretz2011}.

For Bonferroni-based graphical procedures, consonance and sequential rejection
play an important role. In particular, if the monotonicity condition
\[
w_j(J) \leq w_j(J')
\qquad \text{for all } J' \subseteq J \subseteq I
\text{ and } j \in J'
\]
is satisfied, then the resulting closed testing procedure is consonant
\citep{Hommel2007,Bretz2011}. This leads to a sequentially rejective shortcut
procedure, so that the graphical multiple test can be carried out through
iterative graph updates rather than explicit evaluation of all
$2^m - 1$ intersection hypotheses.

\subsection{Group sequential extension}

The graphical framework has been extended to repeated testing over interim and
final analyses by \citet{MaurerBretz2013}. Following their notation, consider
testing $h > 1$ one-sided hypotheses $H_i$, $i \in I = \{1,\ldots,h\}$, in a
group sequential trial with $k$ analysis time points $t = 1,\ldots,k$.

Let $I_t$ denote the amount of statistical information available at time point
$t$, and let $y_t = I_t / I_{\max}$ denote the corresponding information
fraction, where $I_{\max}$ is the planned maximum information if early stopping
does not occur. A family of spending functions is a nondecreasing function
$a(\gamma,y)$ in the information time $y$, parameterized by the overall
significance level $\gamma$, such that
\[
a(\gamma,0)=0
\qquad \text{and} \qquad
a(\gamma,1)=\gamma.
\]

Let $Z_t$ denote the test statistic for a single hypothesis at time point $t$,
and let $p_t$ denote the corresponding nominal $p$-value. The nominal
significance levels $\alpha_t^*(\gamma)$ are defined through
\[
\alpha_t(\gamma)
=
a(\gamma,y_t) - a(\gamma,y_{t-1})
=
P\left(
\{p_t \leq \alpha_t^*(\gamma)\}
\cap
\bigcap_{s=1}^{t-1}
\{p_s > \alpha_s^*(\gamma)\}
\right),
\]
so that the cumulative probability of rejection up to time point $t$ is
controlled by the chosen spending function \citep{LanDeMets1983,MaurerBretz2013}.

For repeated testing of multiple hypotheses, Maurer and Bretz extend the
Bonferroni-based graphical framework by applying weighted group sequential
Bonferroni tests to the intersection hypotheses. For each $J \subseteq I$, the
weights $w_i(J)$ determine the local significance levels for the hypotheses in
$H_J$, and the corresponding nominal rejection boundaries at look $t$ are
obtained from the relevant spending functions with $\gamma$ replaced by
$w_i(J)\alpha$. Under mild monotonicity conditions on the spending functions,
the resulting procedure remains sequentially rejective, allowing the graphical
approach to be used in group sequential trials essentially in the same way as
in the fixed-design case \citep{MaurerBretz2013}.



\section{Availability and reproducibility} \label{sec:availability}
The \pkg{graphMTP} application is available at
\url{https://oncotrialdesign.shinyapps.io/graphMTP/}.
Source code is available at
\url{https://github.com/IDSWG-Oncology/graphMTP/}.

\bibliography{references}

\end{document}